% =====================================================================
%  Banco complementar --- Divisor de Tensao  (REVISADO)
%  Substitui a versao gerada por template (8 clones em N2, 9 em N3 e
%  8 questoes conceituais marcadas como N4).
%  O cap04 ja cobre: maior parcela (MC), tres resistores em serie,
%  100 V com 60/40, tres iguais sob 60 V, projeto de 5 V a partir de
%  15 V, saida em vazio e com carga R_L, potenciometro de 10 kohm como
%  divisor, malha unica com cinco resistores, divisor para carga de
%  10 mA, divisor 4k/2k carregado e sistema de 5 V para carga de
%  100 ohm. O EFEITO DE CARGA basico, portanto, ja esta no capitulo --
%  e tambem no banco de circuitos mistos. Aqui ele so reaparece na
%  forma da CONDICAO DE PROJETO (N4) e da nao linearidade do
%  potenciometro carregado, que sao angulos novos.
%  Distribuicao mantida: 4 N1 + 8 N2 + 9 N3 + 8 N4 = 29
% =====================================================================

\subsubsection*{Banco complementar --- Divisor de tensão}

\begin{atencao}[Organização do banco]
A relação $U_o=U_{in}\dfrac{R_2}{R_1+R_2}$ pressupõe \textbf{corrente comum} ---
ou seja, saída em vazio. Todo o resto do tópico decorre de levar a sério essa
hipótese: o divisor tem resistência interna $R_1\parallel R_2$, e qualquer carga
comparável a ela derruba a saída. O N3 explora tomadas múltiplas, sensores e
fontes não ideais; o N4 trata de projeto sob restrição de potência, tolerância
de razão, sonda compensada e a comparação com reguladores.
\end{atencao}

\blocofixacao

\begin{questao}[1]{}
Três resistores de $\SI{1}{\kilo\ohm}$, $\SI{2}{\kilo\ohm}$ e
$\SI{3}{\kilo\ohm}$ estão em série sob $\SI{12}{\volt}$. Determine a tensão
sobre o resistor \emph{intermediário}.
\resposta{$U=\SI{4}{\volt}$}
\resolucao{A regra do divisor vale para qualquer número de elementos em série:
cada um fica com a fração de sua própria resistência sobre o total.
\[
U_2=12\cdot\frac{2}{1+2+3}=12\cdot\frac13=\SI{4}{\volt}.
\]
\textbf{Conferência:} as três tensões são $2$, $4$ e $\SI{6}{\volt}$, somando
$\SI{12}{\volt}$ \checkmark, e guardam a proporção $1{:}2{:}3$ dos
resistores.}
\end{questao}

\begin{questao}[1]{}
Um divisor deve entregar $25\%$ da tensão de entrada. Determine a razão
$R_1/R_2$.
\resposta{$R_1/R_2=3$}
\resolucao{
\[
\frac{R_2}{R_1+R_2}=0{,}25
\;\Longrightarrow\;
R_1+R_2=4R_2
\;\Longrightarrow\;
R_1=3R_2.
\]
\textbf{Observação importante:} apenas a \emph{razão} é fixada --- $3\,$k$/1\,$k
e $30\,$k$/10\,$k produzem a mesma saída em vazio. O que os distingue é a
corrente de repouso e a resistência interna do divisor, e é isso que decide o
projeto quando há carga.}
\end{questao}

\begin{questao}[1]{}
Em um divisor em vazio, ambos os resistores são \textbf{dobrados}. A tensão de
saída:
\begin{alternativas}[2]
\item dobra
\item cai à metade
\item não se altera
\item depende dos valores originais
\end{alternativas}
\resposta{(c)}
\resolucao{A razão $\dfrac{2R_2}{2R_1+2R_2}=\dfrac{R_2}{R_1+R_2}$ é
inalterada --- a saída em vazio depende só da proporção.\\
\textbf{O que muda:} a corrente de repouso cai à metade (menos consumo) e a
resistência interna $R_1\parallel R_2$ \emph{dobra} (mais sensível à carga).
Escalar um divisor troca consumo por robustez, e essa é sempre a decisão
central do projeto.}
\end{questao}

\begin{questao}[1]{}
Analise os casos-limite de um divisor: (a) $R_2\to0$; (b) $R_2\to\infty$.
\resposta{(a) $U_o\to0$; (b) $U_o\to U_{in}$}
\resolucao{(a) Com $R_2=0$, a saída está em \textbf{curto} com o terra:
$U_o=0$. A corrente vale $U_{in}/R_1$ --- toda a potência vai para $R_1$.\\
(b) Com $R_2\to\infty$ (circuito aberto), não circula corrente, não há queda em
$R_1$ e $U_o\to U_{in}$.\\
\textbf{Conclusão:} a saída de um divisor está sempre entre $0$ e $U_{in}$ --- ele
só \emph{reduz} tensão, nunca a eleva. Obter tensão maior que a entrada exige
elemento armazenador (indutor ou capacitor) e comutação.}
\end{questao}

\blocoaplicacao

\begin{questao}[2]{}
Um divisor de $\SI{12}{\volt}$ usa três resistores em série de
$\SI{1}{\kilo\ohm}$, $\SI{1}{\kilo\ohm}$ e $\SI{2}{\kilo\ohm}$, com duas
tomadas intermediárias.
\begin{itensq}
\item Determine as duas tensões de saída em relação ao terra.
\item Determine a corrente de repouso.
\end{itensq}
\resposta{(a) $\SI{9}{\volt}$ e $\SI{6}{\volt}$; (b) $\SI{3}{\milli\ampere}$}
\resolucao{Total: $R=\SI{4}{\kilo\ohm}$ e
$I=12/4000=\SI{3}{\milli\ampere}$.\\
A tomada mais alta tem, abaixo de si, $1+2=\SI{3}{\kilo\ohm}$:
\[
U_A=12\cdot\frac{3}{4}=\SI{9}{\volt}.
\]
A tomada inferior tem $\SI{2}{\kilo\ohm}$ abaixo:
\[
U_B=12\cdot\frac{2}{4}=\SI{6}{\volt}.
\]
\textbf{Cuidado:} cada saída conta apenas a resistência \emph{entre ela e o
terra}, não a do trecho acima. E as duas tomadas se \emph{influenciam}: carregar
uma altera a outra, pois compartilham a mesma cadeia.}
\end{questao}

\begin{questao}[2]{}
Uma fonte de $\SI{30}{\volt}$ alimenta dois resistores iguais em série, e o
ponto médio é adotado como referência (terra). Determine as tensões dos dois
extremos.
\resposta{$\SI{+15}{\volt}$ e $\SI{-15}{\volt}$}
\resolucao{Cada resistor cai $\SI{15}{\volt}$. Tomando o ponto médio como
$\SI{0}{\volt}$, o terminal superior fica em $\SI{+15}{\volt}$ e o inferior em
$\SI{-15}{\volt}$.\\
\textbf{Aplicação:} é assim que se obtém alimentação \emph{simétrica} para
amplificadores operacionais a partir de uma fonte simples. \\
\textbf{Limitação séria:} o "terra" assim criado é um ponto de divisor, com
resistência interna $R/2$. Qualquer corrente drenada de forma assimétrica pelos
dois ramos desloca a referência. Na prática, o divisor precisa ser reforçado por
um seguidor de tensão ou por capacitores de grande valor.}
\end{questao}

\begin{questao}[2]{}
Um divisor de $\SI{24}{\volt}$ tem $R_1=\SI{8}{\kilo\ohm}$ e
$R_2=\SI{4}{\kilo\ohm}$. Determine a saída, a corrente de repouso e a potência
total dissipada.
\resposta{$\SI{8}{\volt}$; $\SI{2}{\milli\ampere}$; $\SI{48}{\milli\watt}$}
\resolucao{
\[
U_o=24\cdot\frac{4}{12}=\SI{8}{\volt};
\qquad
I=\frac{24}{12000}=\SI{2}{\milli\ampere};
\]
\[
P=U_{in}I=24(\pot{2}{-3})=\SI{48}{\milli\watt}.
\]
\textbf{Distribuição:} $P_1=8000(\pot{2}{-3})^{2}=\SI{32}{\milli\watt}$ e
$P_2=\SI{16}{\milli\watt}$ --- proporcionais às resistências, como em qualquer
série.\\
Note que essa potência é consumida \emph{permanentemente}, mesmo sem carga
alguma conectada. Em equipamentos alimentados por bateria, divisores de repouso
são uma fonte silenciosa de descarga.}
\end{questao}

\begin{questao}[2]{}
Um divisor tem $R_1=\SI{2}{\kilo\ohm}$ fixo e $R_2$ ajustável de
$\SI{0}{\ohm}$ a $\SI{8}{\kilo\ohm}$, sob $\SI{10}{\volt}$. Determine a faixa
da tensão de saída.
\resposta{de $\SI{0}{\volt}$ a $\SI{8}{\volt}$}
\resolucao{
\[
R_2=0 \Rightarrow U_o=0;
\qquad
R_2=\SI{8}{\kilo\ohm} \Rightarrow U_o=10\cdot\frac{8}{10}=\SI{8}{\volt}.
\]
\textbf{Não linearidade do ajuste:} no meio do curso
($R_2=\SI{4}{\kilo\ohm}$), a saída é $10\cdot4/6=\SI{6.67}{\volt}$ --- e não
$\SI{4}{\volt}$, como uma escala linear sugeriria. A relação
$U_o(R_2)$ é uma hipérbole, não uma reta.\\
Para obter ajuste \emph{linear}, é preciso usar um potenciômetro (em que a soma
$R_1+R_2$ permanece constante), e não um reostato em série com resistor fixo.}
\end{questao}

\begin{questao}[2]{}
Um divisor deve entregar $\SI{4.5}{\volt}$ com $R_2=\SI{3}{\kilo\ohm}$.
Determine $U_{in}$ para $R_1=\SI{5}{\kilo\ohm}$.
\resposta{$U_{in}=\SI{12}{\volt}$}
\resolucao{
\[
4{,}5=U_{in}\cdot\frac{3}{5+3}
\;\Longrightarrow\;
U_{in}=4{,}5\cdot\frac{8}{3}=\SI{12}{\volt}.
\]
\textbf{Verificação:} $12\cdot3/8=\SI{4.5}{\volt}$ \checkmark, e a corrente é
$12/8000=\SI{1.5}{\milli\ampere}$.\\
A relação do divisor pode ser resolvida em qualquer das quatro variáveis ---
$U_{in}$, $U_o$, $R_1$ ou $R_2$ --- desde que as outras três sejam conhecidas.}
\end{questao}

\begin{questao}[2]{}
Um termistor NTC vale $\SI{10}{\kilo\ohm}$ a $\SI{25}{\celsius}$ e
$\SI{4}{\kilo\ohm}$ a $\SI{50}{\celsius}$. Ele ocupa a posição de $R_2$ em um
divisor com $R_1=\SI{10}{\kilo\ohm}$ sob $\SI{5}{\volt}$. Determine a saída nas
duas temperaturas.
\resposta{$\SI{2.5}{\volt}$ e $\SI{1.43}{\volt}$}
\resolucao{
\[
\SI{25}{\celsius}:\ U_o=5\cdot\frac{10}{20}=\SI{2.5}{\volt};
\qquad
\SI{50}{\celsius}:\ U_o=5\cdot\frac{4}{14}=\SI{1.43}{\volt}.
\]
Variação de $\SI{1.07}{\volt}$ em $\SI{25}{\celsius}$, ou cerca de
$\SI{43}{\milli\volt\per\celsius}$ na média --- sinal amplo e fácil de digitalizar.\\
\textbf{Ressalva:} a relação \emph{não} é linear, pois nem o termistor é linear
em $T$, nem o divisor é linear em $R_2$. Para converter tensão em temperatura é
preciso tabela, polinômio ou a equação de Steinhart--Hart. A escolha de
$R_1=\SI{10}{\kilo\ohm}$ não é casual --- ela maximiza a sensibilidade em torno
de $\SI{25}{\celsius}$, como se demonstra no bloco de desafio.}
\end{questao}

\begin{questao}[2]{}
Um LDR varia de $\SI{1}{\kilo\ohm}$ (claro) a $\SI{100}{\kilo\ohm}$ (escuro).
Ele é colocado como $R_1$ de um divisor com $R_2=\SI{10}{\kilo\ohm}$ sob
$\SI{5}{\volt}$. Determine a faixa de saída e o sentido da variação.
\resposta{de $\SI{4.55}{\volt}$ (claro) a $\SI{0.45}{\volt}$ (escuro)}
\resolucao{
\[
\text{claro}:\ U_o=5\cdot\frac{10}{11}=\SI{4.55}{\volt};
\qquad
\text{escuro}:\ U_o=5\cdot\frac{10}{110}=\SI{0.45}{\volt}.
\]
Com o LDR na posição \emph{superior}, mais luz significa \emph{maior} tensão de
saída.\\
\textbf{Inversão do comportamento:} trocando as posições (LDR embaixo), a saída
passaria de $\SI{0.45}{\volt}$ no claro a $\SI{4.55}{\volt}$ no escuro. A
escolha da posição define a polaridade do sensor --- decisão de projeto que
custa nada e evita inverter o sinal depois em software.}
\end{questao}

\begin{questao}[2]{}
Um divisor projetado com $R_1=\SI{8.7}{\kilo\ohm}$ e $R_2=\SI{3.3}{\kilo\ohm}$
para $12\to\SI{3.3}{\volt}$ deve usar valores E24. Adotando
$R_1=\SI{9.1}{\kilo\ohm}$, determine a saída real e o erro.
\resposta{$U_o=\SI{3.19}{\volt}$; erro de $-3{,}2\%$}
\resolucao{
\[
U_o=12\cdot\frac{3{,}3}{9{,}1+3{,}3}=12\cdot\frac{3{,}3}{12{,}4}=\SI{3.19}{\volt};
\qquad
\text{erro}=\frac{3{,}19-3{,}30}{3{,}30}=-3{,}2\%.
\]
\textbf{Alternativa melhor:} o par $\SI{10}{\kilo\ohm}/\SI{3.9}{\kilo\ohm}$ dá
$U_o=\SI{3.37}{\volt}$, erro de apenas $+2{,}0\%$. Vale sempre testar
\emph{vários} pares E24 antes de escolher --- a razão importa mais que os
valores individuais.\\
\textbf{Se $3\%$ for demais:} use a série E96 (tolerância $1\%$), que contém
$\SI{8.66}{\kilo\ohm}$ e resolve o problema diretamente.}
\end{questao}

\blocoaprofundamento

\begin{questao}[3]{}
Um divisor $R_1=\SI{6}{\kilo\ohm}$, $R_2=\SI{3}{\kilo\ohm}$ é alimentado por
uma fonte \textbf{real} de $\SI{12}{\volt}$ com $r=\SI{1}{\kilo\ohm}$.
\begin{itensq}
\item Determine a saída, considerando $r$.
\item Compare com a previsão que ignora $r$.
\item Determine a resistência de Thévenin vista na saída.
\end{itensq}
\resposta{(a) $\SI{3.6}{\volt}$; (b) previsão de $\SI{4}{\volt}$, erro de
$10\%$; (c) $R_{Th}=\SI{2.1}{\kilo\ohm}$}
\resolucao{(a) A resistência interna soma-se a $R_1$:
\[
U_o=12\cdot\frac{3}{1+6+3}=12\cdot\frac{3}{10}=\SI{3.6}{\volt}.
\]
(b) Ignorando $r$: $12\cdot3/9=\SI{4}{\volt}$. O erro é de $10\%$ --- nada
desprezível.\\
(c) Zerando a f.e.m. (mas mantendo $r$), a saída enxerga $R_2$ em paralelo com
$(r+R_1)$:
\[
R_{Th}=\frac{3\cdot7}{10}=\SI{2.1}{\kilo\ohm}.
\]
\textbf{Leitura:} a resistência interna da fonte se comporta como parte de
$R_1$ --- ela piora a razão \emph{e} aumenta a impedância de saída. Em divisores
de baixa impedância (resistores de dezenas de ohms), $r$ pode dominar
completamente o projeto.}
\end{questao}

\begin{questao}[3]{}
Dois divisores idênticos ($R_1=\SI{3}{\kilo\ohm}$, $R_2=\SI{1}{\kilo\ohm}$)
são alimentados pela mesma fonte de $\SI{8}{\volt}$, mas em um deles $R_2$ é
trocado por $\SI{1.1}{\kilo\ohm}$.
\begin{itensq}
\item Determine as duas saídas e a diferença entre elas.
\item Determine a sensibilidade da diferença a variações de $R_2$.
\end{itensq}
\resposta{(a) $\SI{2}{\volt}$ e $\SI{2.146}{\volt}$; diferença de
$\SI{146}{\milli\volt}$}
\resolucao{(a)
\[
U_A=8\cdot\frac{1}{4}=\SI{2}{\volt};
\qquad
U_B=8\cdot\frac{1{,}1}{4{,}1}=\SI{2.146}{\volt};
\qquad
\Delta U=\SI{146}{\milli\volt}.
\]
(b) Derivando $U=U_{in}R_2/(R_1+R_2)$:
\[
\frac{\mathrm{d}U}{\mathrm{d}R_2}
=\frac{U_{in}R_1}{(R_1+R_2)^{2}}
=\frac{8(3000)}{(4000)^{2}}=\SI{1.5}{\milli\volt\per\ohm}.
\]
Para $\Delta R_2=\SI{100}{\ohm}$, prevê-se $\SI{150}{\milli\volt}$ --- muito
próximo dos $\SI{146}{\milli\volt}$ exatos, pois a variação é pequena.\\
\textbf{Esta é a estrutura da ponte:} dois divisores comparados, em que o
\emph{desequilíbrio} carrega a informação. A grande vantagem é que variações
comuns aos dois ramos --- deriva da fonte, temperatura ambiente --- se cancelam
na diferença.}
\end{questao}

\begin{questao}[3]{}
Um divisor $R_1=R_2=\SI{100}{\kilo\ohm}$ sob $\SI{10}{\volt}$ é medido por um
voltímetro cuja resistência de entrada é $\SI{1}{\mega\ohm}$.
\begin{itensq}
\item Determine a leitura do instrumento.
\item Determine o erro percentual.
\item Determine a resistência mínima do voltímetro para erro inferior a
      $1\%$.
\end{itensq}
\resposta{(a) $\SI{4.76}{\volt}$; (b) $-4{,}8\%$; (c)
$\ge\SI{5}{\mega\ohm}$}
\resolucao{(a) O voltímetro entra em paralelo com $R_2$:
\[
R_2'=\frac{100\cdot1000}{1100}=\SI{90.9}{\kilo\ohm};
\qquad
U=10\cdot\frac{90{,}9}{100+90{,}9}=\SI{4.76}{\volt}.
\]
(b) Erro: $(4{,}76-5{,}00)/5{,}00=-4{,}8\%$.\\
(c) A resistência interna do divisor é
$R_{Th}=100\parallel100=\SI{50}{\kilo\ohm}$. O erro relativo é
aproximadamente $R_{Th}/R_V$, então
\[
\frac{50\,\mathrm{k}}{R_V}\le0{,}01
\;\Longrightarrow\;
R_V\ge\SI{5}{\mega\ohm}.
\]
\textbf{Lição de instrumentação:} o instrumento \emph{perturba} o circuito que
mede. Multímetros digitais comuns têm $\SI{10}{\mega\ohm}$ --- suficientes aqui,
mas insuficientes para divisores de megaohms. E note que a especificação
depende de $R_{Th}$ do circuito, não da tensão medida.}
\end{questao}

\begin{questao}[3]{}
Um divisor $R_1=\SI{4}{\kilo\ohm}$, $R_2=\SI{1}{\kilo\ohm}$ sob
$\SI{10}{\volt}$ fornece referência a um conversor A/D de $10$ bits com fundo
de escala de $\SI{2.5}{\volt}$.
\begin{itensq}
\item Determine a saída e a margem em relação ao fundo de escala.
\item Determine o degrau de quantização.
\item Compare o erro de quantização com o erro de resistores de $1\%$.
\end{itensq}
\resposta{(a) $\SI{2}{\volt}$, $80\%$ do fundo; (b) $\SI{2.44}{\milli\volt}$;
(c) resistores dominam ($\approx\SI{16}{\milli\volt}$)}
\resolucao{(a) $U_o=10\cdot1/5=\SI{2}{\volt}$, ou $80\%$ do fundo de escala ---
boa utilização da faixa, sem risco de saturação.\\
(b) Com $10$ bits, $2^{10}=1024$ degraus:
\[
\Delta=\frac{2{,}5}{1024}=\SI{2.44}{\milli\volt}.
\]
(c) Com resistores de $1\%$, o pior erro de razão vale
$\dfrac{R_1}{R_1+R_2}(t_1+t_2)=0{,}8(2\%)=1{,}6\%$, ou
$0{,}016\times2=\SI{32}{\milli\volt}$ --- \textbf{treze vezes} o degrau de
quantização.\\
\textbf{Conclusão de projeto:} não adianta usar um conversor de $12$ ou
$16$ bits com esse divisor. O gargalo é o divisor, e a solução é resistores de
$0{,}1\%$, calibração por software, ou uma referência de tensão dedicada.
Aumentar a resolução do conversor sem corrigir o divisor é gastar dinheiro em
dígitos sem significado.}
\end{questao}

\begin{questao}[3]{}
Um divisor de $\SI{48}{\volt}$ para $\SI{12}{\volt}$ deve alimentar uma carga
de $\SI{100}{\ohm}$.
\begin{itensq}
\item Determine $R_1$ e $R_2$ ignorando a carga, com corrente de repouso de
      $\SI{120}{\milli\ampere}$.
\item Determine a saída real com a carga conectada.
\item Avalie a viabilidade.
\end{itensq}
\resposta{(a) $R_1=\SI{300}{\ohm}$, $R_2=\SI{100}{\ohm}$;
(b) $\SI{6.86}{\volt}$; (c) inviável}
\resolucao{(a) $R_2=12/0{,}120=\SI{100}{\ohm}$;
$R_1=(48-12)/0{,}120=\SI{300}{\ohm}$.\\
(b) A carga de $\SI{100}{\ohm}$ fica em paralelo com $R_2$:
\[
R_2'=\frac{100\cdot100}{200}=\SI{50}{\ohm};
\qquad
U_o=48\cdot\frac{50}{350}=\SI{6.86}{\volt}.
\]
Queda de $43\%$ em relação ao projetado.\\
(c) \textbf{Inviável.} A carga consome $\SI{120}{\milli\ampere}$ a
$\SI{12}{\volt}$ --- exatamente a corrente de repouso, e não uma fração dela. A
regra de bolso exige repouso de $10$ a $20$ vezes a corrente da carga, o que
levaria $R_2$ a $\SI{10}{\ohm}$ e a dissipação total a
$48\times1{,}2=\SI{57.6}{\watt}$ para entregar $\SI{1.44}{\watt}$ --- rendimento
de $2{,}5\%$.\\
\textbf{Diagnóstico geral:} divisores resistivos servem para gerar
\emph{referências} de alta impedância, não para \emph{alimentar} cargas. Aqui o
caso pede regulador.}
\end{questao}

\begin{questao}[3]{}
Um divisor com $R_1=\SI{10}{\kilo\ohm}$ e $R_2=\SI{10}{\kilo\ohm}$ opera sob
$\SI{200}{\volt}$.
\begin{itensq}
\item Determine a potência de cada resistor e a tensão sobre cada um.
\item Especifique os resistores considerando potência \emph{e} tensão máxima de
      trabalho.
\end{itensq}
\resposta{$\SI{1}{\watt}$ e $\SI{100}{\volt}$ em cada; exige resistor de
potência com tensão de trabalho adequada}
\resolucao{(a) $I=200/20000=\SI{10}{\milli\ampere}$;
\[
U_1=U_2=\SI{100}{\volt};
\qquad
P_1=P_2=10^{4}(10^{-2})^{2}=\SI{1}{\watt}.
\]
(b) \textbf{Dois critérios independentes:}
\begin{itemize}
\item \textbf{Potência:} com fator $2$ de folga, resistores de
$\SI{2}{\watt}$.
\item \textbf{Tensão máxima de trabalho:} um resistor de filme comum de
$\frac14\,\si{\watt}$ suporta cerca de $\SI{250}{\volt}$; os de
$\SI{2}{\watt}$ chegam a $\SI{500}{\volt}$. Os $\SI{100}{\volt}$ exigidos são
atendidos com folga.
\end{itemize}
\textbf{Por que o segundo critério existe:} a tensão máxima é limitada pela
ruptura do revestimento e pelo arco entre as espiras do filme --- fenômeno que
independe da potência. Em divisores de \emph{alta} tensão (quilovolts), é comum
que a tensão, e não a potência, force o uso de vários resistores em série,
mesmo com dissipação irrisória.}
\end{questao}

\begin{questao}[3]{}
Um divisor deve fornecer $\SI{5}{\volt}$ a partir de $\SI{15}{\volt}$ para uma
entrada de microcontrolador com corrente de entrada de
$\SI{1}{\micro\ampere}$.
\begin{itensq}
\item Determine $R_1$ e $R_2$ para corrente de repouso de
      $\SI{10}{\micro\ampere}$.
\item Verifique o erro introduzido pela corrente de entrada.
\item Comente o compromisso com o consumo.
\end{itensq}
\resposta{(a) $R_1=\SI{1}{\mega\ohm}$, $R_2=\SI{500}{\kilo\ohm}$;
(b) erro de $\approx3\%$}
\resolucao{(a) $R_2=5/\pot{10}{-6}=\SI{500}{\kilo\ohm}$;
$R_1=10/\pot{10}{-6}=\SI{1}{\mega\ohm}$.\\
(b) A corrente de entrada de $\SI{1}{\micro\ampere}$ é $10\%$ da corrente de
repouso --- longe de desprezível. Modelando a entrada como resistência
equivalente $R_{in}=5/\pot{1}{-6}=\SI{5}{\mega\ohm}$:
\[
R_2'=\frac{500\cdot5000}{5500}=\SI{454.5}{\kilo\ohm};
\qquad
U_o=15\cdot\frac{454{,}5}{1454{,}5}=\SI{4.69}{\volt},
\]
erro de $-6{,}2\%$.\\
(c) \textbf{Compromisso.} Reduzir o erro exige aumentar a corrente de repouso ---
com $\SI{100}{\micro\ampere}$ ($R_1=\SI{100}{\kilo\ohm}$,
$R_2=\SI{50}{\kilo\ohm}$), o erro cai para cerca de $0{,}7\%$, mas o consumo em
repouso sobe dez vezes.\\
Em um dispositivo alimentado por bateria de $\SI{1000}{\milli\ampere\hour}$,
$\SI{100}{\micro\ampere}$ contínuos consomem a bateria em pouco mais de um ano
\emph{só com o divisor}. É o tipo de detalhe que decide a autonomia de um
produto.}
\end{questao}

\begin{questao}[3]{}
Determine, em um divisor, a expressão da resistência de Thévenin vista na saída
e avalie-a para $R_1=\SI{6}{\kilo\ohm}$ e $R_2=\SI{3}{\kilo\ohm}$.
\begin{itensq}
\item Deduza $R_{Th}$.
\item Mostre que $R_{Th}<\min(R_1,R_2)$.
\item Explique seu papel no carregamento.
\end{itensq}
\resposta{$R_{Th}=R_1\parallel R_2=\SI{2}{\kilo\ohm}$}
\resolucao{(a) Zerando a fonte de entrada (curto), a saída enxerga $R_1$ e
$R_2$ em paralelo:
\[
R_{Th}=\frac{R_1R_2}{R_1+R_2}=\frac{6\cdot3}{9}=\SI{2}{\kilo\ohm}.
\]
(b) O paralelo de dois resistores positivos é sempre menor que o menor deles ---
resultado geral já estabelecido para associações. Aqui,
$\SI{2}{\kilo\ohm}<\SI{3}{\kilo\ohm}$ \checkmark.\\
(c) \textbf{Papel no carregamento.} Conectada uma carga $R_L$, o divisor se
comporta como fonte de Thévenin:
\[
U_L=U_{Th}\cdot\frac{R_L}{R_L+R_{Th}}.
\]
O erro relativo é aproximadamente $R_{Th}/R_L$ para $R_L\gg R_{Th}$. Logo:
\begin{itemize}
\item $R_L=10R_{Th}$ $\Rightarrow$ erro $\approx9\%$;
\item $R_L=100R_{Th}$ $\Rightarrow$ erro $\approx1\%$.
\end{itemize}
\textbf{Consequência de projeto:} o único parâmetro que importa para o
carregamento é $R_{Th}$ --- não $R_1$, não $R_2$, não a razão. Reduzir os dois
resistores proporcionalmente reduz $R_{Th}$ e melhora a robustez, ao custo de
consumo.}
\end{questao}

\begin{questao}[3]{}
Um divisor alimenta uma carga que \textbf{varia} entre $\SI{10}{\kilo\ohm}$ e
$\SI{100}{\kilo\ohm}$. O divisor tem $R_{Th}=\SI{2}{\kilo\ohm}$ e
$U_{Th}=\SI{5}{\volt}$.
\begin{itensq}
\item Determine a faixa de tensão de saída.
\item Determine a regulação percentual.
\item Indique como melhorá-la.
\end{itensq}
\resposta{(a) de $\SI{4.17}{\volt}$ a $\SI{4.90}{\volt}$;
(b) $\approx15\%$}
\resolucao{(a)
\[
R_L=\SI{10}{\kilo\ohm}:\ U=5\cdot\frac{10}{12}=\SI{4.17}{\volt};
\qquad
R_L=\SI{100}{\kilo\ohm}:\ U=5\cdot\frac{100}{102}=\SI{4.90}{\volt}.
\]
(b) Regulação:
\[
\frac{4{,}90-4{,}17}{4{,}90}=15{,}0\%.
\]
(c) \textbf{Como melhorar.}
\begin{itemize}
\item \textbf{Reduzir $R_{Th}$.} Com $R_{Th}=\SI{200}{\ohm}$, a faixa passa a
$4{,}90$--$\SI{4.99}{\volt}$ e a regulação cai para $1{,}8\%$ --- ao custo de
dez vezes mais consumo de repouso.
\item \textbf{Inserir um seguidor de tensão} (amp-op em configuração de ganho
unitário). A impedância de saída cai para frações de ohm e a regulação
praticamente desaparece, sem aumentar o consumo do divisor.
\item \textbf{Usar regulador}, se a carga também exigir corrente apreciável.
\end{itemize}
\textbf{Observação:} a segunda opção é a mais elegante --- ela desacopla
\emph{razão} de \emph{impedância}, que no divisor puro estão inevitavelmente
amarradas.}
\end{questao}

\blocodesafio

\begin{questao}[4]{}
\textbf{(Condição de carregamento.)} Deduza a condição para que uma carga
$R_L$ altere a saída de um divisor em menos de um erro relativo $\varepsilon$.
\begin{itensq}
\item Obtenha a expressão exata do erro.
\item Deduza a condição aproximada para $\varepsilon$ pequeno.
\item Construa a tabela de $R_L/R_{Th}$ para
      $\varepsilon=10\%$, $1\%$ e $0{,}1\%$.
\end{itensq}
\resposta{$\varepsilon=\dfrac{R_{Th}}{R_L+R_{Th}}$; aproximadamente
$R_L\ge R_{Th}/\varepsilon$}
\resolucao{(a) Com carga,
$U_L=U_{Th}\dfrac{R_L}{R_L+R_{Th}}$, logo
\[
\varepsilon=\frac{U_{Th}-U_L}{U_{Th}}=1-\frac{R_L}{R_L+R_{Th}}
=\frac{R_{Th}}{R_L+R_{Th}}.
\]
(b) Para $R_L\gg R_{Th}$, o denominador tende a $R_L$:
\[
\varepsilon\approx\frac{R_{Th}}{R_L}
\;\Longrightarrow\;
R_L\ge\frac{R_{Th}}{\varepsilon}.
\]
(c)
\begin{center}
\begin{tabular}{ccc}
\hline
$\varepsilon$ & $R_L/R_{Th}$ (exato) & $R_L/R_{Th}$ (aprox.)\\
\hline
$10\%$ & $9$ & $10$\\
$1\%$ & $99$ & $100$\\
$0{,}1\%$ & $999$ & $1000$\\
\hline
\end{tabular}
\end{center}
A aproximação erra por exatamente $1$ unidade e é conservadora --- pode ser
usada sem receio.\\
\textbf{Regra de bolso resultante:} para $1\%$ de erro,
$R_L\ge100\,R_{Th}$. Note que $R_{Th}=R_1\parallel R_2$ é sempre \emph{menor}
que qualquer um dos resistores, o que torna a condição menos severa do que
comparar $R_L$ com $R_2$ diretamente --- erro comum que leva a
superdimensionar o divisor.}
\end{questao}

\begin{questao}[4]{}
\textbf{(Projeto sob restrição de potência.)} Projete um divisor que reduza
$\SI{100}{\volt}$ a $\SI{5}{\volt}$, dissipando no máximo
$\SI{0.5}{\watt}$.
\begin{itensq}
\item Determine $R_1$ e $R_2$.
\item Determine a potência e a tensão de cada resistor.
\item Determine a maior carga que pode ser atendida com erro de até $1\%$.
\end{itensq}
\resposta{(a) $R_1=\SI{19}{\kilo\ohm}$, $R_2=\SI{1}{\kilo\ohm}$;
(b) $0{,}475$ e $\SI{0.025}{\watt}$; (c) $R_L\ge\SI{95}{\kilo\ohm}$}
\resolucao{(a) A restrição de potência fixa a \emph{soma}:
\[
P=\frac{U_{in}^{2}}{R_1+R_2}\le0{,}5
\;\Longrightarrow\;
R_1+R_2\ge\frac{100^{2}}{0{,}5}=\SI{20}{\kilo\ohm}.
\]
A razão fixa a \emph{proporção}: $R_2/(R_1+R_2)=0{,}05$. Adotando o mínimo,
\[
R_2=\SI{1}{\kilo\ohm};\qquad R_1=\SI{19}{\kilo\ohm}.
\]
(b) $I=100/20000=\SI{5}{\milli\ampere}$;
\[
P_1=19000(\pot{5}{-3})^{2}=\SI{0.475}{\watt};
\qquad
P_2=\SI{0.025}{\watt};
\]
com $U_1=\SI{95}{\volt}$ e $U_2=\SI{5}{\volt}$.\\
\textbf{Especificação:} $R_1$ concentra $95\%$ da dissipação --- use
$\SI{1}{\watt}$ nele e $\frac18\,\si{\watt}$ em $R_2$. Alternativamente,
dividir $R_1$ em dois de $\SI{9.5}{\kilo\ohm}$ reparte o calor e a tensão.\\
(c) $R_{Th}=19\parallel1=\SI{950}{\ohm}$; pela condição anterior,
\[
R_L\ge100\,R_{Th}=\SI{95}{\kilo\ohm}.
\]
\textbf{Tensão de projeto:} o divisor entrega $\SI{5}{\volt}$ a uma carga que
consome no máximo $\SI{53}{\micro\ampere}$ --- serve para um conversor A/D, não
para alimentar nada.}
\end{questao}

\begin{questao}[4]{}
\textbf{(Potenciômetro carregado.)} Um potenciômetro de $\SI{10}{\kilo\ohm}$
é usado como divisor sob $\SI{10}{\volt}$, com carga
$R_L=\SI{10}{\kilo\ohm}$ na saída. Seja $x$ a fração do curso.
\begin{itensq}
\item Deduza $U_o(x)$ com a carga.
\item Avalie em $x=0{,}25$, $0{,}5$ e $0{,}75$ e compare com o caso ideal.
\item Determine onde o desvio é máximo e como reduzi-lo.
\end{itensq}
\resposta{Em $x=0{,}5$: $\SI{4.00}{\volt}$ contra $\SI{5}{\volt}$ ideal ---
desvio de $20\%$}
\resolucao{(a) A porção inferior vale $xR_p$ e fica em paralelo com $R_L$; a
superior vale $(1-x)R_p$:
\[
U_o=U_{in}\,\frac{xR_p\parallel R_L}{(1-x)R_p+\left(xR_p\parallel R_L\right)}.
\]
Com $R_p=R_L=\SI{10}{\kilo\ohm}$, tem-se
$xR_p\parallel R_L=\dfrac{10x}{1+x}\,\si{\kilo\ohm}$ e
\[
U_o=10\cdot\frac{\dfrac{10x}{1+x}}{10(1-x)+\dfrac{10x}{1+x}}
=10\cdot\frac{x}{(1-x)(1+x)+x}
=\frac{10x}{1+x-x^{2}}.
\]
(b)
\begin{center}
\begin{tabular}{cccc}
\hline
$x$ & ideal & com carga & desvio\\
\hline
$0{,}25$ & $\SI{2.50}{\volt}$ & $\SI{2.105}{\volt}$ & $-15{,}8\%$\\
$0{,}50$ & $\SI{5.00}{\volt}$ & $\SI{4.000}{\volt}$ & $-20{,}0\%$\\
$0{,}75$ & $\SI{7.50}{\volt}$ & $\SI{6.316}{\volt}$ & $-15{,}8\%$\\
\hline
\end{tabular}
\end{center}
(c) O desvio \emph{absoluto} é máximo perto do meio do curso, onde a
resistência de Thévenin do potenciômetro ($x(1-x)R_p$) atinge seu máximo,
$R_p/4=\SI{2.5}{\kilo\ohm}$. Nos extremos ($x\to0$ ou $x\to1$) a impedância
tende a zero e a carga deixa de importar.\\
\textbf{A curva deixa de ser reta:} o potenciômetro linear, carregado, comporta-se
de forma \textbf{não linear}, e o erro depende da posição --- o que é
inaceitável em um controle graduado.\\
\textbf{Correções:} \emph{(i)} usar $R_L\ge100R_p$, o que reduz o desvio máximo
a menos de $0{,}25\%$; \emph{(ii)} inserir um seguidor de tensão entre o cursor
e a carga, solução padrão em instrumentação; \emph{(iii)} usar potenciômetro de
lei "logarítmica reversa" pré-distorcida --- gambiarra que só funciona para uma
carga específica.}
\end{questao}

\begin{questao}[4]{}
\textbf{(Tolerância de razão.)} Um divisor usa resistores de $1\%$ com entrada
exata.
\begin{itensq}
\item Deduza a expressão do pior erro relativo da saída.
\item Avalie para as razões $12\to\SI{6}{\volt}$, $12\to\SI{5}{\volt}$ e
      $100\to\SI{5}{\volt}$.
\item Explique por que resistores \emph{casados} melhoram drasticamente o
      resultado.
\end{itensq}
\resposta{$\varepsilon=\dfrac{R_1}{R_1+R_2}(t_1+t_2)$; $1{,}00\%$,
$1{,}17\%$ e $1{,}90\%$}
\resolucao{(a) De $U_o=U_{in}\dfrac{R_2}{R_1+R_2}$, tomando logaritmos e
diferenciando:
\[
\frac{\Delta U_o}{U_o}
=\frac{R_1}{R_1+R_2}\left(\frac{\Delta R_2}{R_2}-\frac{\Delta R_1}{R_1}\right),
\]
cujo módulo máximo (desvios opostos) é
$\varepsilon=\dfrac{R_1}{R_1+R_2}(t_1+t_2)$.\\
(b) O fator $\dfrac{R_1}{R_1+R_2}$ vale $1-U_o/U_{in}$:
\begin{center}
\begin{tabular}{ccc}
\hline
Razão & $R_1/(R_1+R_2)$ & Pior erro\\
\hline
$12\to6$ & $0{,}500$ & $1{,}00\%$\\
$12\to5$ & $0{,}583$ & $1{,}17\%$\\
$100\to5$ & $0{,}950$ & $1{,}90\%$\\
\hline
\end{tabular}
\end{center}
\textbf{Padrão:} quanto \emph{maior} a razão de divisão, pior o erro --- ele
tende a $t_1+t_2=2\%$ quando $U_o\ll U_{in}$, e nunca o ultrapassa.\\
(c) \textbf{Resistores casados.} Em uma rede resistiva integrada, os dois
resistores são fabricados lado a lado, no mesmo processo e no mesmo substrato.
Seus desvios são \textbf{correlacionados}: se um sai $0{,}8\%$ acima do
nominal, o outro também sai. O termo entre parênteses,
$\dfrac{\Delta R_2}{R_2}-\dfrac{\Delta R_1}{R_1}$, tende a \emph{zero}.\\
Na prática, redes casadas oferecem razão com precisão de $0{,}05\%$ mesmo com
valores absolutos de $\pm20\%$. \textbf{O que importa no divisor é a razão, não
os valores} --- e o casamento é a única forma barata de obtê-la com precisão.
Esse é o princípio que viabiliza conversores A/D e amplificadores de
instrumentação.}
\end{questao}

\begin{questao}[4]{}
\textbf{(Sonda atenuadora.)} Uma sonda $10{:}1$ é formada por
$\SI{9}{\mega\ohm}$ na ponta, em série com a entrada de
$\SI{1}{\mega\ohm}$ do osciloscópio.
\begin{itensq}
\item Verifique a razão de atenuação e a resistência de entrada total.
\item Determine o erro ao medir uma fonte com $R_{Th}=\SI{100}{\kilo\ohm}$,
      com e sem a sonda.
\item Explique o que se perde ao usar a sonda.
\end{itensq}
\resposta{(a) $1/10$ e $\SI{10}{\mega\ohm}$; (b) $1{,}0\%$ com sonda,
$9{,}1\%$ sem}
\resolucao{(a) A entrada do osciloscópio \emph{é} o $R_2$ do divisor:
\[
\frac{U_o}{U_{in}}=\frac{1}{9+1}=\frac{1}{10}\ \checkmark;
\qquad
R_{entrada}=9+1=\SI{10}{\mega\ohm}.
\]
(b) \textbf{Com sonda:} o instrumento apresenta $\SI{10}{\mega\ohm}$ à fonte:
\[
\varepsilon=\frac{100\,\mathrm{k}}{10\,\mathrm{M}+100\,\mathrm{k}}=1{,}0\%.
\]
\textbf{Sem sonda} (cabo direto, $\SI{1}{\mega\ohm}$):
\[
\varepsilon=\frac{100\,\mathrm{k}}{1\,\mathrm{M}+100\,\mathrm{k}}=9{,}1\%.
\]
A sonda reduz o erro de carregamento por um fator $9$.\\
(c) \textbf{O que se perde: amplitude.} O sinal chega ao instrumento dez vezes
menor, o que reduz a relação sinal-ruído em $\SI{20}{\decibel}$ e torna
inviável medir sinais de poucos milivolts.\\
\textbf{O compromisso central da sonda} é exatamente esse: troca-se sensibilidade
por não perturbação. Por isso os osciloscópios trazem seletor $1\times/10\times$
--- sinais pequenos em fontes de baixa impedância pedem $1\times$; sinais grandes
em circuitos de alta impedância pedem $10\times$.\\
\emph{(Em sinais variáveis há ainda a capacitância parasita, que exige o
capacitor de compensação ajustável presente em toda sonda real --- assunto de
corrente alternada.)}}
\end{questao}

\begin{questao}[4]{}
\textbf{(Otimização de sensor.)} Um termistor de resistência $R_t$ ocupa a
posição inferior de um divisor com resistor fixo $R$ e alimentação $U$.
\begin{itensq}
\item Obtenha $\dfrac{\mathrm{d}U_o}{\mathrm{d}R_t}$.
\item Determine o valor de $R$ que maximiza a sensibilidade em um dado ponto de
      operação.
\item Avalie para $U=\SI{5}{\volt}$, $R_t=\SI{10}{\kilo\ohm}$ e coeficiente de
      $\SI{-400}{\ohm\per\celsius}$.
\end{itensq}
\resposta{(b) $R=R_t$; (c) $\SI{-50}{\milli\volt\per\celsius}$}
\resolucao{(a) De $U_o=U\dfrac{R_t}{R+R_t}$:
\[
\frac{\mathrm{d}U_o}{\mathrm{d}R_t}
=U\,\frac{(R+R_t)-R_t}{(R+R_t)^{2}}
=\frac{U\,R}{(R+R_t)^{2}}.
\]
(b) Maximizando em relação a $R$, com $R_t$ fixo:
\[
\frac{\mathrm{d}}{\mathrm{d}R}\left[\frac{R}{(R+R_t)^{2}}\right]
=\frac{(R+R_t)-2R}{(R+R_t)^{3}}
=\frac{R_t-R}{(R+R_t)^{3}}=0
\;\Longrightarrow\;
\boxed{R=R_t}
\]
--- e a saída no ponto ótimo é exatamente $U/2$.\\
\textbf{Mesma estrutura} do resultado de máxima potência: a derivada de
$R/(R+R_t)^{2}$ é a mesma que a de $P_R$ em uma série, e por isso o ótimo cai
em igualdade.\\
(c) Com $R=R_t=\SI{10}{\kilo\ohm}$:
\[
\frac{\mathrm{d}U_o}{\mathrm{d}R_t}=\frac{5(10^{4})}{(\pot{2}{4})^{2}}
=\SI{1.25e-4}{\volt\per\ohm},
\]
\[
\frac{\mathrm{d}U_o}{\mathrm{d}T}
=(\pot{1.25}{-4})(-400)=\SI{-50}{\milli\volt\per\celsius}.
\]
\textbf{Excelente sensibilidade:} um conversor A/D de $10$ bits com fundo de
$\SI{5}{\volt}$ tem degrau de $\SI{4.9}{\milli\volt}$, resolvendo cerca de
$\SI{0.1}{\celsius}$.\\
\textbf{Ressalva:} o ótimo vale para \emph{um} ponto de operação. Em faixa
ampla de temperatura, $R_t$ varia por ordens de grandeza e a sensibilidade cai
nas extremidades --- escolha $R$ igual a $R_t$ na temperatura de maior interesse,
não na média da faixa.}
\end{questao}

\begin{questao}[4]{}
\textbf{(Divisor contra regulador.)} Compare as duas soluções para obter
$\SI{5}{\volt}$ a partir de $\SI{12}{\volt}$, para uma carga que varia de
$\SI{1}{\milli\ampere}$ a $\SI{10}{\milli\ampere}$.
\begin{itensq}
\item Dimensione o divisor com repouso de dez vezes a carga máxima e determine
      a regulação.
\item Determine o rendimento de cada solução.
\item Estabeleça o critério de escolha.
\end{itensq}
\resposta{(a) $R_1=\SI{70}{\ohm}$, $R_2=\SI{50}{\ohm}$; regulação
$\approx5{,}3\%$; (b) $3{,}8\%$ contra $42\%$}
\resolucao{(a) Repouso de $\SI{100}{\milli\ampere}$:
$R_2=5/0{,}1=\SI{50}{\ohm}$ e
$R_1=7/0{,}1=\SI{70}{\ohm}$.
\[
R_{Th}=50\parallel70=\SI{29.2}{\ohm};
\qquad
\Delta U=R_{Th}\,\Delta I_L=29{,}2(\pot{9}{-3})=\SI{0.26}{\volt},
\]
ou seja, cerca de $5{,}3\%$ de variação.\\
(b) \textbf{Divisor:} consumo total
$\approx12\times0{,}11=\SI{1.32}{\watt}$ para entregar no máximo
$5\times0{,}01=\SI{50}{\milli\watt}$ --- rendimento de $3{,}8\%$.\\
\textbf{Regulador linear:} consome a mesma corrente da carga,
$12\times0{,}01=\SI{120}{\milli\watt}$ para entregar
$\SI{50}{\milli\watt}$ --- rendimento de $41{,}7\%$ (igual à razão
$U_o/U_{in}$), com regulação típica melhor que $0{,}1\%$.\\
(c) \textbf{Critério.}
\begin{itemize}
\item \textbf{Divisor} apenas quando a carga é de \emph{altíssima impedância} e
praticamente constante: entradas de conversores A/D, polarização de bases,
referências internas. Vantagens: dois componentes, custo desprezível, sem ruído
de comutação.
\item \textbf{Regulador} sempre que a carga consumir corrente apreciável ou
variar. O regulador linear resolve com rendimento $U_o/U_{in}$; o chaveado
ultrapassa $90\%$, ao custo de complexidade e ruído.
\end{itemize}
\textbf{Fronteira prática:} se a corrente da carga for maior que cerca de
$1\%$ da corrente de repouso que o divisor exigiria, o divisor já não é a
escolha certa.}
\end{questao}

\begin{questao}[4]{}
\textbf{(Divisor com proteção.)} Projete a interface entre um sinal de entrada
que varia de $\SI{0}{\volt}$ a $\SI{30}{\volt}$ e um conversor A/D de
$\SI{3.3}{\volt}$ e impedância de entrada muito alta.
\begin{itensq}
\item Determine a razão e escolha os resistores.
\item Verifique o comportamento em sobretensão de $\SI{50}{\volt}$ e proponha
      proteção.
\item Determine o consumo e discuta o compromisso com a impedância exigida pelo
      conversor.
\end{itensq}
\resposta{(a) razão $\approx1/10$: $R_1=\SI{91}{\kilo\ohm}$,
$R_2=\SI{10}{\kilo\ohm}$; (c) $\SI{297}{\micro\ampere}$ no fundo de escala}
\resolucao{(a) É preciso $30\to\SI{3.3}{\volt}$, razão $0{,}11$. Com
$R_2=\SI{10}{\kilo\ohm}$ e $R_1=\SI{91}{\kilo\ohm}$ (ambos E24):
\[
U_o=30\cdot\frac{10}{101}=\SI{2.97}{\volt},
\]
$10\%$ abaixo do fundo de escala --- margem prudente, que acomoda tolerância dos
resistores sem saturar o conversor.\\
(b) Com $\SI{50}{\volt}$ na entrada, a saída iria a
$50\times10/101=\SI{4.95}{\volt}$, acima do limite absoluto do conversor.\\
\textbf{Proteção:} um diodo Zener de $\SI{3.3}{\volt}$ (ou um diodo de
grampeamento para a alimentação) em paralelo com $R_2$. Em sobretensão, ele
conduz e limita a saída; a corrente é confinada por $R_1$ a
\[
I=\frac{50-3{,}3}{91\,\mathrm{k}}=\SI{0.51}{\milli\ampere},
\]
valor inofensivo tanto para o Zener quanto para $R_1$
($P=\SI{24}{\milli\watt}$). \textbf{É $R_1$ que faz a proteção funcionar} ---
sem ele, o diodo teria de absorver corrente ilimitada.\\
(c) No fundo de escala, $I=30/101\,\mathrm{k}=\SI{297}{\micro\ampere}$ e
$P\approx\SI{8.9}{\milli\watt}$.\\
\textbf{Compromisso.} $R_{Th}=91\parallel10=\SI{9}{\kilo\ohm}$. Conversores A/D
de aproximações sucessivas exigem impedância de fonte baixa (tipicamente
$<\SI{10}{\kilo\ohm}$) para carregar o capacitor de amostragem no tempo
disponível --- este projeto fica no limite.\\
\textbf{Soluções:} reduzir os resistores por um fator $10$ (custa dez vezes mais
consumo), acrescentar um capacitor de $\SI{100}{\nano\farad}$ na entrada do
conversor (fornece a carga instantânea de amostragem), ou inserir um seguidor
de tensão. A segunda opção é a mais comum, e a mais barata.}
\end{questao}
